\section{引言}

Mamba通过输入相关的状态更新实现线性序列建模，但扫描前仍然需要确定Token顺序 \cite{gu2023mamba}。Vision Mamba通常采用栅格、方向、窗口或其他预先设计的路线 \cite{zhu2024vim,liu2024vmamba,yang2024plainmamba,huang2024localmamba}。这些路线便于并行计算，却不一定符合当前图像自身的结构。

本文的直觉很简单。人观察图像时，不会对所有区域投入相同精度：宽泛背景可以先获得粗略认识，醒目的局部细节则得到更晚、更密集的读取。本文不声称模拟生物视觉，只把这一想法用于安排选择性状态空间模型的信息顺序。低显著性信息较早进入全局历史，高显著性信息较晚进入；热力变化快的位置获得更密的采样。每次Mamba扫描开始前，完整顺序都已经确定。

\method{}从显著性标量场生成这个顺序。闭合等值线分量形成局部路线，等高线树记录这些分量随显著性升高而出现、延续、分裂、汇合和终止的关系 \cite{carr2003contourtrees}。因此不需要把每个圈强制匹配到最近的下一层圈，也不会在盆地周围产生错误的远距离连接。细小树枝可以保留在拓扑中，但不一定消耗扫描Token。

初版使用传统算法生成显著性场，并直接从原始RGB图像读取视觉内容，不包含CNN特征Stem。同一张Sobel场统一完成三个决定：带符号的法线响应为整圈选择低显著性读取侧，响应强度决定圈内采样密度，最强的平滑响应位置确定公共锚点 $p_0$。这样不再为方向、密度和起止点分别设计互不相关的规则。

局部读取与全局汇总严格分开。每个保留轮廓独立编码：Normal Mamba从碰撞点读回带Tag的轮廓点；两条Contour Mamba从同一个 $p_0$ 出发，按相反方向绕圈，并在带Tag的 $p_0$ 副本结束。每条树枝的终止圈读取两侧法线，再进行第二次闭环汇集。所有局部编码完成后，Tree Mamba才沿稀疏增广等高线树从低到高处理轮廓摘要。图~\ref{fig:local}展示按照这些规则生成的一条完整路线。

本文贡献可概括为：
\begin{itemize}
    \item 用闭合等值线与低到高的增广等高线树构造随输入变化的视觉序列；
    \item 用一套Sobel规则统一决定整圈读取侧、固定预算内的采样密度和公共锚点；
    \item 用带Tag的原图像素法线与闭环编码器完成局部读取，再以Tree Mamba处理分裂、无序汇合、终止叶和空路线。
\end{itemize}

本文给出完整的架构与算法规范。其贡献是确定性路线构建、Token接口、拓扑汇总规则和结构分析；本文不作实证性能声明。
